%%%%%%%%%%%%%%%%%%%%%%%%%%%%%%%%%%%%%%%%%%%%%%%
\section{ Transformation rules.}
%\label{sec:gauge_rules}
\label{app:gauge_rules}

For a given real gauge function $\chi(\vb{r},t)$, the potentials transform
according to Eqs.~\eqref{eq:gauge-rules-A} and \eqref{eq:gauge-rules},
and each Kohn--Sham orbital acquires a position-dependent phase,
\begin{equation}
  \Psi_j'(\vb{r},t) = \hat U_\chi\,\Psi_j(\vb{r},t),
  \qquad
  \hat U_\chi(\vb{r},t) = \exp\!\left[\frac{\imath q_e}{\hbar c}\,\chi(\vb{r},t)\right],
  \qquad
  \hat U_\chi^\dagger\hat U_\chi = \hat I .
  \label{eq:orbital_rule}
\end{equation}
Since $\hat U_\chi$ acts by multiplication, the transformation of any operator
$\hat O$ is most transparent in terms of its position-space kernel,
\begin{equation}
  (\hat O\Psi)(\vb{r}) = \int O(\vb{r},\vb{r}')\,\Psi(\vb{r}')\,d^3r' .
  \nonumber
\end{equation}
The transformed operator $\hat O' = \hat U_\chi\hat O\,\hat U_\chi^\dagger$ has
the kernel
\begin{equation}
  O'(\vb{r},\vb{r}';t)
  = e^{\imath q_e\chi(\vb{r},t)/\hbar c}\;O(\vb{r},\vb{r}')\;
    e^{-\imath q_e\chi(\vb{r}',t)/\hbar c}.
  \label{eq:kernel_rule}
\end{equation}
Equation~\eqref{eq:kernel_rule} applies to every operator; its consequences
depend only on the structure of the kernel, which distinguishes three classes
of operators summarized in Table~\ref{tab:gauge_operator_classes}.

\emph{(i) Multiplicative operators.} For an operator diagonal in the position
representation, $O(\vb{r},\vb{r}') = f(\vb{r})\,\delta(\vb{r}-\vb{r}')$, the two
phases in Eq.~\eqref{eq:kernel_rule} are evaluated at the same point and cancel,
\begin{equation}
  \hat U_\chi f(\vb{r})\,\hat U_\chi^\dagger = f(\vb{r}),
  \qquad
  |\Psi_j'(\vb{r},t)|^2 = |\Psi_j(\vb{r},t)|^2 .
  \label{eq:local_rule}
\end{equation}
The position operator, the electron density, the external scalar potential
$q_e\Phi$, and the local ionic, Hartree, and adiabatic exchange--correlation
potentials are therefore unchanged.

\emph{(ii) Differential operators.} For kernels containing derivatives of
$\delta(\vb{r}-\vb{r}')$, the derivative also acts on the phase. Applying the
product rule to an arbitrary function gives, for the canonical momentum
$\vb{\hat p}=-\imath\hbar\nabla$,
\begin{equation}
  \hat U_\chi\,\vb{\hat p}\,\hat U_\chi^\dagger
  = \vb{\hat p} - \frac{q_e}{c}\nabla\chi ,
  \label{eq:p_rule}
\end{equation}
i.e., every gradient is replaced by $\nabla-(\imath q_e/\hbar c)\nabla\chi$.
The kinetic (mechanical) momentum
$\hat{\boldsymbol{\pi}}[\vb{A}] = \vb{\hat p} - (q_e/c)\vb{A}$
therefore transforms covariantly,
\begin{equation}
  \hat U_\chi\,\hat{\boldsymbol{\pi}}[\vb{A}]\,\hat U_\chi^\dagger
  = \vb{\hat p} - \frac{q_e}{c}\left(\vb{A}+\nabla\chi\right)
  = \hat{\boldsymbol{\pi}}[\vb{A}'] ,
  \label{eq:pi_rule}
\end{equation}
and the kinetic energy retains its minimal-coupling form,
$\hat U_\chi\hat{\boldsymbol{\pi}}^2\hat U_\chi^\dagger = \hat{\boldsymbol{\pi}}'^{\,2}$.
For a general phase $e^{\imath\beta\chi}$, Eq.~\eqref{eq:p_rule} would read
$\vb{\hat p}-\hbar\beta\nabla\chi$; requiring this minimal-coupling form with
$\vb{A}'$ fixes $\beta = q_e/(\hbar c)$, the phase used in
Eq.~\eqref{eq:orbital_rule}. The canonical momentum itself retains its
coordinate representation $-\imath\hbar\nabla$ in every gauge; by
Eq.~\eqref{eq:p_rule} it is not the transform of itself, and its expectation
value is gauge dependent.

The time derivative in the TDKS equation behaves in the same way with respect
to time, $\hat U_\chi\,\imath\hbar\partial_t\,\hat U_\chi^\dagger
= \imath\hbar\partial_t + (q_e/c)\,\partial_t\chi$. Because the Hamiltonian
generates the time evolution, this produces an additional term: the
transformed orbitals obey $\imath\hbar\,\partial_t\Psi_j' = \hat H'\Psi_j'$ with
\begin{equation}
  \hat H' = \hat U_\chi\hat H\hat U_\chi^\dagger
          + \imath\hbar\,(\partial_t\hat U_\chi)\hat U_\chi^\dagger
          = \hat U_\chi\hat H\hat U_\chi^\dagger - \frac{q_e}{c}\,\frac{\partial\chi}{\partial t}.
  \label{eq:H_rule}
\end{equation}
The last term combines with $q_e\Phi$ to give $q_e\Phi'$. The two potential
rules, Eqs.~\eqref{eq:gauge-rules-A} and \eqref{eq:gauge-rules}, are thus the
spatial and temporal manifestations of the same noncommutativity of the phase
with differential operators.

\emph{(iii) Integral (nonlocal) operators.} For a general kernel
$V(\vb{r},\vb{r}')$ the phases are evaluated at different points and do not
cancel,
\begin{equation}
  V'(\vb{r},\vb{r}';t)
  = \exp\!\left\{\frac{\imath q_e}{\hbar c}\big[\chi(\vb{r},t)-\chi(\vb{r}',t)\big]\right\}
    V(\vb{r},\vb{r}') .
  \label{eq:nonlocal_rule}
\end{equation}
Here it is useful to distinguish fixed external operators from
orbital-dependent ones. Fixed nonlocal operators, such as the separable projectors of
norm-conserving pseudopotentials and the projectors used in DFT+$U$, must be
transformed explicitly according to Eq.~\eqref{eq:nonlocal_rule}; transforming
the kinetic term alone breaks gauge
invariance.\cite{IsmailBeigi-2001-EM,pickard-mauri-2001,Mattiat-Luber-2022}
Orbital-dependent nonlocal operators, such as Fock exchange in hybrid
functionals, are built from the transformed orbitals and acquire the phase of
Eq.~\eqref{eq:nonlocal_rule} automatically.

\emph{Observables.} A time-independent observable transforms according to
Eq.~\eqref{eq:kernel_rule},
\begin{equation}
  \hat O' = \hat U_\chi\,\hat O\,\hat U_\chi^\dagger,
  \qquad
  \langle\Psi_j'|\hat O'|\Psi_j'\rangle = \langle\Psi_j|\hat O|\Psi_j\rangle ,
  \label{eq:observable_rule}
\end{equation}
so that the density, $\langle\hat{\boldsymbol{\pi}}\rangle$, and the current
density are gauge invariant. Equations~\eqref{eq:orbital_rule}--\eqref{eq:observable_rule}
are general; they do not yet assume the particular gauge function of
Eq.~\eqref{eq:chi_solution}.

\begin{table*}[tb]
\caption{Behavior of operator classes under the gauge transformation
$\hat O\to\hat U_\chi\hat O\,\hat U_\chi^\dagger$,
$\hat U_\chi=\exp[\imath q_e\chi(\vb{r},t)/\hbar c]$. All rules follow from the
kernel transformation, Eq.~\eqref{eq:kernel_rule}.}
\label{tab:gauge_operator_classes}
\centering
\small
\setlength{\tabcolsep}{4pt}
\renewcommand{\arraystretch}{1.3}
% No extra packages needed: each cell starts with \raggedright and rows end
% with \tabularnewline (plain \\ is redefined by \raggedright).
\begin{tabular}{p{0.175\textwidth}p{0.15\textwidth}p{0.18\textwidth}p{0.16\textwidth}p{0.16\textwidth}}
\hline
\raggedright Class
 & \raggedright Kernel $O(\vb{r},\vb{r}')$
 & \raggedright Rule
 & \raggedright Examples
 & \raggedright Consequence \tabularnewline
\hline
\raggedright Multiplicative
 & \raggedright $f(\vb{r})\,\delta(\vb{r}-\vb{r}')$
 & \raggedright unchanged, Eq.~\eqref{eq:local_rule}
 & \raggedright $\vb{r}$, $\rho$, $q_e\Phi$, $V^{\mathrm{loc}}_{\mathrm{ion}}$, $V_{\mathrm H}$, $V_{\mathrm{xc}}$
 & \raggedright invariant \tabularnewline
\raggedright Differential (space)
 & \raggedright derivatives of $\delta(\vb{r}-\vb{r}')$
 & \raggedright $\vb{\hat p}\to\vb{\hat p}-(q_e/c)\nabla\chi$, Eq.~\eqref{eq:p_rule}
 & \raggedright $\vb{\hat p}$, $\hat{\boldsymbol{\pi}}$, $\hat T$
 & \raggedright $\vb{A}\to\vb{A}+\nabla\chi$ \tabularnewline
\raggedright Differential (time)
 & \raggedright ---
 & \raggedright $\imath\hbar\partial_t\to\imath\hbar\partial_t+(q_e/c)\partial_t\chi$
 & \raggedright TDKS time derivative
 & \raggedright $\Phi\to\Phi-c^{-1}\partial_t\chi$, Eq.~\eqref{eq:H_rule} \tabularnewline
\raggedright Integral (nonlocal)
 & \raggedright general $V(\vb{r},\vb{r}')$
 & \raggedright two-point phase, Eq.~\eqref{eq:nonlocal_rule}
 & \raggedright $\hat V_{\mathrm{NL}}$, DFT+$U$ projectors, Fock exchange
 & \raggedright $[\hat{\vb{r}},\hat V]$ terms in $\hat H$ and current \tabularnewline
\hline
\end{tabular}
\end{table*}
